\begin{table}[!htbp]
\centering
\footnotesize
\setlength{\tabcolsep}{3.5pt}
\caption*{\textbf{Table C14.} Assumptions and allocations that permit a non-positive slope}
\begin{tabularx}{\linewidth}{lCCCC}
\toprule
Election & Out-of-zone voters within (points) & Largest additional tolerance s with a positive slope & First reversing allocation: movers minus stayers, conservative share (points) & Fewest residents' early ballots that must depart to reverse \\
\midrule
2022 & 10 & none & 12 to 21 & 16,685 \\
2022 & 20 & none & 14 to 18 & 15,340 \\
2024 & 10 & 0.021 & 11 to 17 & 32,177 \\
2024 & 20 & none & 9 to 14 & 30,103 \\
2025 & 10 & none & 20 to 24 & 13,539 \\
2025 & 20 & none & 21 to 23 & 12,512 \\
2026 & 10 & $\infty $ & no reversal & no reversal \\
2026 & 20 & 0.087 & $-$2 to 24 & 69,384 \\
\bottomrule
\end{tabularx}
\par\vspace{2pt}{\raggedright\footnotesize\textit{Notes.} Mobility-calibrated movement. s: residents' early ballots may depart from their neighborhood benchmark (election-day share shifted by the district's early-vote gap) by their station's departure plus s; $\infty $: positive however much they depart; none: not positive even when they depart only as much as their station. First reversing allocation: the one with the lowest slope at s just above that value; movers voted at another station of their zone, stayers at their own; range over price terciles. Last column: among allocations with a non-positive slope, the fewest residents' conservative early ballots departing from the benchmark (out-of-zone departures and concentration in Online Appendix Table C16). Each reversal is minimal only in its own metric, with other assumptions fixed.\par}
\end{table}
